\documentclass[11pt,a4paper]{ctexart}
\usepackage[margin=21mm]{geometry}
\usepackage{amsmath,amssymb,xcolor}
\setlength{\parindent}{0pt}
\setlength{\parskip}{5pt}
\newcommand{\R}{\mathbb R}
\newcommand{\T}{\mathbb T}
\newcommand{\dd}{\,\mathrm d}
\newcommand{\sep}{\par\vspace{5pt}\noindent\rule{\linewidth}{0.7pt}\par\vspace{5pt}}
\newcommand{\red}[1]{{\color{red}#1}}
\allowdisplaybreaks
\begin{document}
\begin{center}\Large 小作文1\end{center}
给定 $\T^2$ 上的光滑向量场 $X$，考虑构造光滑向量场 $Y,Z$，使得 $[Y,Z]=X$。
将 $\T^2$ 写成 $\R^2/(2\pi\mathbb Z)^2$。通过覆盖映射将 $X$ 提升到 $\R^2$ 后，
问题等价于在 $\R^2$ 上求解 $[Y,Z]=X$，并要求 $Y,Z$ 对所有
$(2k\pi,2p\pi)$，$k,p\in\mathbb Z$，均具有周期性。
以下仍用 $X$ 表示提升后的向量场，取 $\R^2$ 上坐标为 $(s,t)$。

Lie 括号满足
\[\red{[fU,gV]=fg[U,V]+fU(g)V-gV(f)U.}\]
\sep
\textbf{初步构造与周期性限制。}
先考虑 $[\partial_t,W]=X$。令
\[X=F\partial_s+G\partial_t,\qquad W=a(s,t)\partial_s+b(s,t)\partial_t.\]
则 $[\partial_t,W]=a_t\partial_s+b_t\partial_t$，故方程化为
\[a_t=F,\qquad b_t=G.\]
在 $\R^2$ 上，可直接积分求解。但若要求 $a,b$ 关于 $t$ 有周期性，
则相应右端的周期积分必须为零；特别地，必须有
\[\int_0^{2\pi}F(s,t)\dd t=0\qquad\forall s.\]
因此这一构造不能直接处理一般的 $X$。
\sep
\textbf{旋转公式。} 设 $U,V$ 为光滑向量场，$\theta$ 为光滑函数，令
\[Y=\cos\theta\,U+\sin\theta\,V,\qquad Z=-\sin\theta\,U+\cos\theta\,V.\]
即对两个向量场作角度为 $\theta$ 的旋转。由括号公式，
\begin{align*}
[Y,Z]&=[\cos\theta\,U+\sin\theta\,V,-\sin\theta\,U+\cos\theta\,V]\\
&=\bigl(-\cos\theta\,U(\sin\theta)+\sin\theta\,U(\cos\theta)\bigr)U\\
&\quad+\bigl(\sin\theta\,V(\cos\theta)-\cos\theta\,V(\sin\theta)\bigr)V\\
&\quad+\cos^2\theta[U,V]+\cos\theta\,U(\cos\theta)V-\cos\theta\,V(\cos\theta)U\\
&\quad+\sin^2\theta[U,V]+\sin\theta\,V(-\sin\theta)U+\sin\theta\,U(\sin\theta)V\\
&=[U,V]+\bigl(-\cos^2\theta\,U(\theta)-\sin^2\theta\,U(\theta)\bigr)U\\
&\quad+\cos\theta\sin\theta\,V(\theta)U-\sin\theta\cos\theta\,V(\theta)U\\
&\quad+\bigl(-\sin^2\theta\,V(\theta)-\cos^2\theta\,V(\theta)
             -\cos\theta\sin\theta\,U(\theta)+\sin\theta\cos\theta\,U(\theta)\bigr)V\\
&=[U,V]-U(\theta)U-V(\theta)V.
\end{align*}
取 $U=\partial_t$，$V=W$，则
\begin{align*}
[Y,Z]&=a_t\partial_s+b_t\partial_t-\theta_t\partial_t
 -(a\theta_s+b\theta_t)(a\partial_s+b\partial_t)\\
&=(a_t-a^2\theta_s-ab\theta_t)\partial_s
 +(b_t-ab\theta_s-(1+b^2)\theta_t)\partial_t.
\end{align*}
所以需要求解
\[\begin{cases}F=a_t-a^2\theta_s-ab\theta_t,\\
G=b_t-ab\theta_s-(1+b^2)\theta_t.\end{cases}\]
\sep
\textbf{角度函数的两种选取。}
若取 $\theta=\theta(t)$，使 $\theta_s=0$，则方程化为
\[\begin{cases}a_t-ab\theta_t=F,\\
b_t-(1+b^2)\theta_t=G.\end{cases}\]
固定 $\theta$ 后，可先由第二条 Riccati 方程求 $b$，再由第一条线性方程求 $a$。
但周期性仍然带来限制。事实上，
\[\theta_t+\frac{G}{1+b^2}=\frac{b_t}{1+b^2}=(\arctan b)_t.\]
若 $b$ 是周期函数，且 $\theta(2\pi)-\theta(0)=2\pi k$，$k\in\mathbb Z$，
则在一个周期上积分得到
\[2\pi k+\int_0^{2\pi}\frac{G}{1+b^2}\dd t=0.\]
特别地，若该积分位于 $(0,\frac32\pi)$，便与这一整数倍条件矛盾。
因此，这一选取会引入额外的周期约束。

另一种选取是 $\theta(s)=s$，此时 $\theta_s=1$，方程化为
\[F=a_t-a^2,\qquad G=b_t-ab.\]
这一方程组中不再出现上一种选取中的 $1+b^2$ 系数。
对第一条在一个周期上积分，得到必要条件
\[\int_0^{2\pi}F(s,t)\dd t=-\int_0^{2\pi}a(s,t)^2\dd t\leq0.\]
\sep
\textbf{负周期积分条件下的构造。}
引入待选正常数 $L$，取 $U=L\partial_t$，$V=W/L$，$\theta=s$，即
\[Y=L\cos s\,\partial_t+\frac{\sin s}{L}W,\qquad
Z=-L\sin s\,\partial_t+\frac{\cos s}{L}W.\]
旋转公式给出
\[[Y,Z]=\partial_tW-\frac{a}{L^2}W.\]
故需要求解
\[\tag{1}\begin{cases}a_t-\dfrac{a^2}{L^2}=F,\\[3pt]
b_t-\dfrac{a}{L^2}b=G.\end{cases}\]
\red{假设对每个 $s$，都有 $\int_0^{2\pi}F(s,t)\dd t<0$。}
下面证明：当 $L$ 充分大时，方程组 (1) 存在光滑周期解。

令
\[m(s)=-\frac1{2\pi}\int_0^{2\pi}F(s,t)\dd t>0,\qquad
r(s,t)=\int_0^t(F(s,\tau)+m(s))\dd\tau.\]
则 $r_t=F+m$，且 $r$ 关于 $s,t$ 均为 $2\pi$ 周期。
取 $K>\sup|r|$，再取 $L$ 充分大，使 $L\min_s\sqrt{m(s)}>2K$。设
\[a_\pm(s,t)=-L\sqrt{m(s)}+r(s,t)\pm K.\]
于是
\[a_-<-L\sqrt{m(s)}<a_+<0,\]
并且
\[(a_+)_t=F+m>F+\frac{a_+^2}{L^2},\qquad
(a_-)_t=F+m<F+\frac{a_-^2}{L^2}.\]
因此 $a_+,a_-$ 分别是严格周期上解、下解。

固定 $s$，记 $a(s,t;\xi)$ 为初值问题
\[a_t=F+\frac{a^2}{L^2},\qquad a(s,0;\xi)=\xi\]
的解。初值在 $[a_-(s,0),a_+(s,0)]$ 内时，解被上下解夹住，
因而在 $[0,2\pi]$ 上存在。定义 Poincar\'e 映射
\[P_s(\xi)=a(s,2\pi;\xi).\]
由严格比较，
\[P_s(a_-(s,0))>a_-(s,2\pi)=a_-(s,0),\]
\[P_s(a_+(s,0))<a_+(s,2\pi)=a_+(s,0).\]
所以连续函数 $\xi\mapsto P_s(\xi)-\xi$ 在两个端点异号，存在
$\xi_s$ 使得 $P_s(\xi_s)=\xi_s$。于是 $a(s,t;\xi_s)$ 是严格负的周期解。

严格负的周期解唯一。若 $a_1,a_2$ 是两个不同的负周期解，由解的唯一性，
其差处处不为零，且
\[(a_1-a_2)_t=\frac{a_1^2-a_2^2}{L^2},\qquad
\bigl(\ln|a_1-a_2|\bigr)_t=\frac{a_1+a_2}{L^2}<0,\]
这与周期性矛盾。

再证明 $\xi_s$ 关于 $s$ 光滑。它是 $P_s(\xi)-\xi=0$ 的根，
只需验证 $\partial_\xi P_s(\xi_s)\ne1$。对初值 $\xi$ 求导，令
$V(t)=\partial_\xi a(s,t;\xi_s)$，则
\[V_t=\frac{2a}{L^2}V,\qquad V(0)=1.\]
因此
\[\partial_\xi P_s(\xi_s)=V(2\pi)
=\exp\!\left(\int_0^{2\pi}\frac{2a}{L^2}\dd t\right)<1.\]
由隐函数定理，$\xi_s$ 光滑，故所得周期解 $a$ 光滑。
由 $F$ 关于 $s$ 的周期性及负周期解的唯一性，$a$ 关于 $s$ 也为 $2\pi$ 周期。
\sep
\textbf{线性方程的周期解。} 给定上述 $a$，考虑 $b_t-\frac{ab}{L^2}=G$。其解为
\[b(s,t)=\exp\!\left(\int_0^t\frac{a(s,u)}{L^2}\dd u\right)
\left[b(s,0)+\int_0^t
e^{-\int_0^\tau\frac{a(s,\rho)}{L^2}\dd\rho}G(s,\tau)\dd\tau\right].\]
周期条件 $b(s,0)=b(s,2\pi)$ 等价于
\begin{align*}
b(s,0)
&=\exp\!\left(\int_0^{2\pi}\frac{a(s,u)}{L^2}\dd u\right)\\[-2pt]
&\quad\times\left[b(s,0)+\int_0^{2\pi}
e^{-\int_0^\tau\frac{a(s,\rho)}{L^2}\dd\rho}G(s,\tau)\dd\tau\right].
\end{align*}
由于 $a<0$，外面的指数因子严格小于 $1$，因此可以唯一解出 $b(s,0)$。
由此得到光滑的周期解 $b$；其关于 $s$ 的周期性也由唯一性得到。
代回 $W,Y,Z$，即完成负周期积分条件下的构造。
\sep
\textbf{去掉负周期积分条件。}
令 $X=P\partial_x+Q\partial_y$，并用横线表示相对于标准面积的平均。

\red{Step 1. 假设 $\overline P=\dfrac1{4\pi^2}\int_{\T^2}P\ne0$，不妨 $\overline P>0$。}
若非零平均分量为 $Q$，可先交换坐标；平均为负时可反转相应坐标。
对正整数 $N$，取环面坐标变换
\[\phi_N(s,t)=(t,s+Nt)=(x,y).\]
其整数矩阵行列式为 $-1$，故为环面微分同胚，并且
\[\mathrm d\phi_N(\partial_s)=\partial_y,\qquad
\mathrm d\phi_N(\partial_t)=\partial_x+N\partial_y.\]
因此拉回向量场为
\[\phi_N^*X=
\bigl(Q(t,s+Nt)-NP(t,s+Nt)\bigr)\partial_s
+P(t,s+Nt)\partial_t.\]
为应用前面的构造，需要选取充分大的 $N$，使对所有 $s$ 均有
\[\int_0^{2\pi}\bigl(Q(t,s+Nt)-NP(t,s+Nt)\bigr)\dd t<0.\]
这由下面的平均引理保证。

\textbf{平均引理。} 对于 $h\in C(\T^2)$，有
\[\frac1{2\pi}\int_0^{2\pi}h(s+Nt,t)\dd t
\longrightarrow\frac1{4\pi^2}\int_{\T^2}h\]
其中 $N$ 沿正整数趋于无穷，且收敛关于 $s$ 一致。

\textbf{证明。} 分段换元并利用周期性，得到
\begin{align*}
\frac1{2\pi}\int_0^{2\pi}h(s+Nt,t)\dd t
&=\frac1{2\pi}\sum_{j=0}^{N-1}\int_0^{2\pi}
\frac1N h\!\left(s+u,\frac{2\pi j+u}{N}\right)\dd u\\
&=\frac1{2\pi}\int_0^{2\pi}\left[\sum_{j=0}^{N-1}
h\!\left(s+u,\frac{2\pi j+u}{N}\right)\frac1N\right]\dd u\\
&\longrightarrow\frac1{2\pi}\int_0^{2\pi}
\left[\frac1{2\pi}\int_0^{2\pi}h(s+u,v)\dd v\right]\dd u\\
&=\frac1{4\pi^2}\int_0^{2\pi}\dd u\int_0^{2\pi}h(s+u,v)\dd v
=\frac1{4\pi^2}\int_{\T^2}h.
\end{align*}

括号中的和是 Riemann 和；由 $h$ 在紧环面上的一致连续性，
该收敛关于 $s,u$ 一致，故上述结论关于 $s$ 一致成立。

将引理应用于交换变量后的 $P$，得到
$\frac1{2\pi}\int_0^{2\pi}P(t,s+Nt)\dd t\to\overline P>0$，且关于 $s$ 一致。
又因为 $Q$ 有界，故充分大的同一个整数 $N$ 使上述负积分条件对所有 $s$ 成立。
对 $\phi_N^*X$ 应用前面的构造，再由 $\phi_N$ 推前，即得 $X$ 的括号表示。
\sep
\red{Step 2. 假设 $\overline P=\overline Q=0$。}
零向量场可直接取 $Y=Z=0$。对 $X\not\equiv0$，使用下面的引理。

\textbf{引理（此处不证明）。} 设 $X\not\equiv0$ 是 $\T^2$ 上的光滑向量场，
且 $\overline X=0$。则存在光滑微分同胚
\[\phi:\T^2\to\T^2,\qquad \overline{\phi_*X}\ne0.\]
由此引理，$\phi_*X$ 满足 Step 1 的条件，故可写成两个光滑向量场的括号。
再通过 $\phi$ 拉回，便得到 $X$ 的括号表示。
\clearpage
\sep
\textbf{$T^2\times M$ 的情形（$M$ 紧致）}

设 $M$ 是紧致光滑流形，$s,t$ 是 $T^2$ 上周期为 $2\pi$ 的角坐标。
利用乘积结构，将 $T^2\times M$ 上的光滑向量场写成
\[X=F\partial_s+G+H\partial_t,\qquad G(s,p,t)\in T_pM.\]
假设对每个 $(s,p)\in S^1\times M$，都有
\[\int_0^{2\pi}F(s,p,t)\dd t<0.\]
则同样可以构造 $T^2\times M$ 上的光滑向量场 $Y,Z$，使得 $[Y,Z]=X$。

\textbf{证明。} 取充分大的正常数 $L$，令
\[W=a\partial_s+B+c\partial_t,\qquad B(s,p,t)\in T_pM,\]
并设
\[
Y=L\cos s\,\partial_t+\frac{\sin s}{L}W,
\qquad
Z=-L\sin s\,\partial_t+\frac{\cos s}{L}W.
\]
由旋转公式，
\[[Y,Z]=\partial_tW-\frac{a}{L^2}W.\]
因此只需要求出下面方程组的光滑周期解：
\[
\begin{cases}
a_t-\dfrac{a^2}{L^2}=F,\\[6pt]
B_t-\dfrac{a}{L^2}B=G,\\[6pt]
c_t-\dfrac{a}{L^2}c=H.
\end{cases}
\]

把 $(s,p)$ 视为参数，沿用前面的上下解、Poincar\'e 映射及隐函数定理论证。
由于 $S^1\times M$ 紧致，可以统一选取充分大的 $L$，得到光滑的严格负周期解 $a$。
后两条是在固定的向量空间 $T_pM$ 及 $\mathbb R$ 中的线性方程；其周期乘子
\[\exp\!\left(\int_0^{2\pi}\frac{a(s,p,t)}{L^2}\dd t\right)\in(0,1),\]
故各有唯一的光滑周期解 $B,c$。代回即得 $[Y,Z]=X$。
这里不需要 $TM$ 平凡。取 $M=T^{n-2}$，即得到 $T^n$ 在上述条件下的结论。
\end{document}
